\paragraph{Sparse strata and gradient variance.}
At the anchor draw, mark stratification audits none of the 61 reference-positive training candidates without any weak mark, and every audited label in this stratum is multiplied by $1/\pi=50$. With the image-only auxiliary, the realized signed residual correction in this stratum ranges from $-36.51$ to approximately zero, compared with $-0.36$ to approximately zero for the discriminative auxiliary. Using the weak marks thus reduces the auxiliary error in this stratum (Tab.~\ref{tab:aux}), although the deployed head remains limited by both ranking and calibration.

For candidates without any weak mark, the diagonal design variance of the discriminative moment is $5.67{\times}10^{-6}$ under base-budget mark allocation and $1.44{\times}10^{-6}$ under half-budget paired-tail allocation. Weighting the gradient by the frozen leverage $h_i$ gives $3.22{\times}10^{-6}$ and $4.10{\times}10^{-9}$, respectively; this diagnostic is computed on the same tail that guided the allocation. Candidates away from the tail can still receive large inverse probabilities: under half-budget paired-tail allocation, the scaled residuals in this stratum range from $-0.91$ to $8.45$. The power-tuned multipliers at the anchor draw are 0.67, 0.90 and 0.38, and over the six base-budget mark draws the power-tuned and untuned discriminative heads yield identical omission and alert counts at the primary threshold. All-target recall, in contrast, is sensitive to the head objective across the development grid of ridge penalties and loss weights.

\paragraph{Controlled bookmark cap.}
In this setting, the first reader retains only its largest mark in each scan, which creates 151 imposed omissions on the test set. Across the three repeated draws, EM augmentation recovers 29--37 of these omissions, compared with 25 for first-read training, 32--34 for IPW and 27--35 for discriminative augmentation. The corresponding alert counts are 4--5 of 105 negative patients for EM augmentation, 4 for first-read training and 4--6 for both IPW and discriminative augmentation. Since the cap imposes a size-based selection of omissions, it poses a different problem from the observed omissions studied in the main paper.

\paragraph{Per-target changes.}
Averaged over the six draws, the selected method newly detects 9.1\% [7.3, 11.1] of all targets relative to base-budget mark IPW and loses 6.1\% [4.8, 7.5], with paired patient-bootstrap intervals; for the omissions, the corresponding fractions are 0.8\% [0.0, 2.1] and 4.8\% [2.4, 7.8]. At the anchor draw, these changes amount to eight gained and nine lost detections over all targets, including no gained and two lost omissions. In the calibration-matched analysis, the mean test alert rates are 1.7\% for the selected method and 2.9\% for mark IPW. When audit draws and calibration patients are resampled jointly (six draws by 300 resamples), the central 95\% of the omission effect spans $-14.3$ to $+7.1$ points, with the test patients held fixed. All comparisons in this section evaluate different thresholds on the same held-out roster, and no independent adjudication establishes whether the disputed contours correspond to in-scope nodules.
